This section summarizes features and functionalities that were considered during the design of the smart plant monitoring system but will not be implemented in the current prototype due to limitations in time, resources, and technical complexity. These features may be included in future versions of the system.


\subsection{Mobile application support}
\subsubsection{Description}
A dedicated mobile application will allow users to monitor plant health, receive notifications, and control system settings remotely.
\subsubsection{Source}
User experience enhancement goals
\subsubsection{Constraints}
Requires additional development time, cross-platform compatibility, and backend integration.
\subsubsection{Standards}
Mobile application development standards (iOS/Android guidelines)
\subsubsection{Priority}
Future

\subsection{Advanced AI-based plant disease detection}
\subsubsection{Description}
The system will use machine learning models to detect specific plant diseases from captured images and provide detailed diagnostics and recommendations.
\subsubsection{Source}
Project design discussion
\subsubsection{Constraints}
Requires large datasets, model training, and higher computational resources.
\subsubsection{Standards}
AI/ML best practices and data handling guidelines
\subsubsection{Priority}
Future

\subsection{Cloud-based data analytics and sharing}
\subsubsection{Description}
The system will support cloud integration for storing plant data, enabling long-term analytics and sharing data across multiple devices.
\subsubsection{Source}
System scalability considerations
\subsubsection{Constraints}
Requires cloud infrastructure, security implementation, and ongoing maintenance.
\subsubsection{Standards}
Cloud computing and data security standards
\subsubsection{Priority}
Future